\scanpage{557}
% archive sheet number: 556
\mspage{468}

supposer que $k$ soit un corps, et aussi en prenant la condition $\mathrm{Hom}_{\unsure{\mathcal M|M}}(A,\unsure{Z|z})=\varnothing$, où $Z$ est un objet quelconque de $\mathcal M$. Ce genre de questions semble \unsure{étrangement|étonnamment} délicat\ldots

4.16.2. Il est possible, plus généralement, que la question suivante ait une réponse affirmative : Soit
\[
f:\mathcal M\unsure{\text{---}|\longrightarrow}\mathcal N^{\circ}
\]
un foncteur accessible, avec $\mathcal M,\unsure{\mathcal N|N}$ \inserted{cat.} accessibles, et soit $\underline T$ une sous catégorie de $\mathcal N$ accessible dans $\mathcal N$, et telle qu'il existe un cardinal $\pi$ ayant les propriétés
suivantes : $\mathcal N^\circ$ \deleted{\unsure{est}} satisfait $L_\pi$, \deleted{et $T^\circ$} (i.e. $\mathcal N$ stable par $\varprojlim$ grdes devant $\pi$)
et $T^\circ$ dans $\mathcal N^\circ$ est $L_\pi$-stable (i.e. $T$ stable dans $\mathcal N$ par $\varprojlim$ grdes \unsure{dev. à|dev.|devant} $\unsure{\pi|\pi_0}$).
Alors $\underline S\overset{\unsure{\mathrm{déf}|\mathrm{def}}}{=}f^{-1}(T)$ est-elle accessible dans $\mathcal M$ ?
% slanted
\margin{Faux, avec $\mathcal N=\mathrm{Ens}$ et un foncteur représentable $\mathcal M\to\mathrm{Ens}^{\circ}$, et prenant $\underline T=\{\varnothing\}$, cf \unsure{annotations marginales|annotation marginale} p. 466 et 467.}
Ceci généralise le cas $\mathcal N=(\mathrm{Ens})$, $\underline T=\{\varnothing\}$, $f$ foncteur représenté par $z\in\mathrm{Ob}\,\mathcal M$, envisagé \deleted{\unsure{tout|plus}} à l'heure. Avec le $=$ \deleted{\unsure{de|d}}type d'exemple, mais en prenant $\mathcal N=(\mathrm{Ens})^{\unsure{a|\alpha|*}}$
